\clearpage\thispagestyle{otFront}\phantomsection\label{references}\pdfbookmark[0]{Sources et édition}{bm-references}\otPageTitre{Sources et édition}{Références officielles · France}\fontsize{11}{15}\selectfont
\par\medskip\textbf{1. Arrêté du 18 février 2026 : calendrier}\par « Programmes d’enseignement de français et de mathématiques du cycle des approfondissements (cycle 4) », arrêté du 18-2-2026, JO du 4-3-2026, NOR MENE2602912A, Bulletin officiel n° 10 du 5 mars 2026. Article 3 : application en classe de troisième à la rentrée de l’année scolaire 2028-2029.\par\href{https://www.education.gouv.fr/sites/default/files/document/Bulletin\%20officiel\%20n\%C2\%B0\%2010\%20du\%205\%20mars\%202026-480719.pdf}{\textcolor{otGreenDark}{Consulter la source officielle}}\par
\par\medskip\textbf{2. Programme de mathématiques du cycle 4}\par Annexe 2 de cet arrêté, « Programme de mathématiques pour le cycle 4 » : principes, puis rubriques de Troisième (automatismes, objectifs d’apprentissage et prolongements possibles).\par\href{https://www.education.gouv.fr/sites/default/files/document/Annexe\%202\%20\%E2\%80\%93\%20Programme\%20de\%20math\%C3\%A9matiques\%20pour\%20le\%20cycle\%204-480716.pdf}{\textcolor{otGreenDark}{Consulter la source officielle}}\par
\par\medskip\textbf{3. Exemples pour la mise en œuvre}\par Éduscol, « Mathématiques. Exemples pour la mise en œuvre des programmes 2026. Cycle 4 » : exemples de réussite et conseils de mise en œuvre pour chaque objectif de Troisième.\par\href{https://eduscol.education.gouv.fr/sites/default/files/document/exemple-mise-en-oeuvre-c4-mathematiques0pdf-126080.pdf}{\textcolor{otGreenDark}{Consulter la source officielle}}\par
\par\medskip\textbf{4. Définition de l’épreuve du brevet}\par Note de service du 11-9-2026, NOR MENE2623228N, « Définition des épreuves de l’examen du diplôme national du brevet à compter de la session 2027 », Bulletin officiel spécial n° 4 du 17 septembre 2026 ; partie 2 : épreuve écrite de mathématiques.\par\href{https://www.education.gouv.fr/sites/default/files/document/bulletin-officiel-special-ndeg-4-du-17-septembre-2026-521113.pdf}{\textcolor{otGreenDark}{Consulter la source officielle}}\par
\par\medskip\textbf{5. Maîtrise de la langue aux examens}\par Note de service du 11-9-2026, NOR MENE2623195N, « Exigences en matière de maîtrise de la langue à tous les examens à compter de la session 2027 », même Bulletin officiel spécial n° 4.\par\href{https://www.education.gouv.fr/sites/default/files/document/bulletin-officiel-special-ndeg-4-du-17-septembre-2026-521113.pdf}{\textcolor{otGreenDark}{Consulter la source officielle}}\par
\par\bigskip Sources consultées le 23 octobre 2026. Le programme et les exemples
de mise en œuvre ont guidé le périmètre et la matrice de correspondance. Les
exercices sont originaux, non reproduits des documents officiels. Les défis et
les sujets du livre ne sont pas des sujets officiels.\par\medskip
À cette date, aucun texte propre à la session 2029 du brevet n’était publié.
La note du 11 septembre 2026 s’applique « à compter de la session 2027 » et
précise que les sujets portent sur le programme de Troisième s’il est en vigueur :
c’est la définition retenue pour les brevets blancs. Un texte ultérieur pourrait
la modifier.

\clearpage
